\subsection{Ingeniería de requisitos}

\subsubsection{Definición y alcance}

La ingeniería de requisitos reúne las actividades que permiten
identificar, analizar, especificar, validar y gestionar las necesidades
y restricciones de un sistema a lo largo de su ciclo de
vida~\cite{ieee29148,pohl2025,swebok2024}.

Una necesidad expresa el problema que un interesado quiere resolver.
Para convertirla en un requisito, el equipo debe aclarar qué
comportamiento espera del sistema y en qué condiciones debe ocurrir.
Registrar la fuente de esa decisión permite revisarla cuando el equipo
encuentra una excepción o cambia el alcance del
proyecto~\cite{ieee29148,pohl2025}.

\textbf{Ejemplo en UteqFit-AI.}
La necesidad de conocer la ocupación del gimnasio se traduce en
RF-06, que pide mostrar cuántos estudiantes tienen una entrada
registrada sin su salida correspondiente. Para revisar ese requisito,
se puede preparar un caso con tres entradas y dos salidas asociadas.
El resultado esperado es que el sistema muestre una persona dentro.

\subsubsection{Actores del proceso}

En ingeniería de requisitos, los interesados son personas u
organizaciones que tienen interés en el sistema o pueden verse
afectadas por él. Sus necesidades y restricciones aportan información
para definir lo que el sistema debe hacer~\cite{ieee29148,pohl2025}.

El ingeniero de requisitos organiza las consultas y documenta las
decisiones. El cliente o responsable del proyecto confirma las
prioridades y los usuarios explican cómo realizan las operaciones
cotidianas. El desarrollador aporta información sobre la viabilidad
técnica y ayuda a aclarar las condiciones necesarias para implementar
y comprobar los requisitos~\cite{pohl2025,swebok2024}.

\textbf{Ejemplo en UteqFit-AI.}
Para revisar RF-06, el equipo del PFC puede preparar los casos de
prueba y consultar al encargado de turno sobre qué sucede cuando
un estudiante se va sin registrar su salida. El responsable del
gimnasio puede decidir cómo debe tratarse ese caso. El desarrollador
puede valorar qué información haría falta para corregir el registro
y conservar el historial de la operación.